\documentclass[11pt,a4paper]{article}
\usepackage[margin=24mm]{geometry}
\usepackage{amsmath,amssymb,amsthm,booktabs,longtable,array,hyperref}
\setlength{\emergencystretch}{2em}
\newtheorem{theorem}{Theorem}
\newtheorem{lemma}[theorem]{Lemma}
\newtheorem{proposition}[theorem]{Proposition}
\newcommand{\egr}{\operatorname{egr}}
\newcommand{\vt}{\mathrm{vt}}
\newcommand{\at}{\mathrm{at}}
\title{Edge-girth-regular graphs: complete generation, numerical bounds,\newline
and the quartic one-hexagon obstruction}
\author{Human Compression project}
\date{18 October 2026}
\begin{document}
\maketitle
\begin{abstract}
We assemble proofs and finite certificates for the principal mathematical
results in the June 2024 edition of Goedgebeur and Jooken's
\emph{Exhaustive generation of edge-girth-regular graphs}. The synthesis
contains a corrected complete-generation theorem, all 142 lower-bound
rows, 84 finite upper witnesses, 47 asserted exact counts, and 118 minima
in two symmetry classes. The printed first-phase generation rule and one
symmetry minimum require explicit corrections. A computer-assisted local
argument excludes quartic graphs on 54 vertices having exactly one hexagon
through each edge; together with the smaller-order argument this gives
$57\le n(4,6,1)\le84$. The companion proofs state the external
classification and classical algebraic premises used in their respective
rows. Neither finite-table completeness nor this new local argument
settles the three remaining assertions of the source's Conjecture 2.
\end{abstract}

\section{Definitions, scope and conclusions}
All graphs are finite, simple and connected unless a partial graph is
explicitly being considered. An $\egr(n,k,g,\lambda)$ graph has $n$
vertices, is $k$-regular, has girth $g$, and has exactly $\lambda>0$
unoriented cycles of length $g$ through every edge. Let $n(k,g,\lambda)$
be its least possible order, with value $\infty$ if no such graph exists.
Write $n_{\vt}$ and $n_{\at}$ for minima restricted to vertex-transitive
and arc-transitive graphs, respectively. These restricted minima need not
be unrestricted minima.

The fixed source throughout is arXiv:2401.08271v2, dated 25 June 2024
\cite{source}. Its physical pages 17--20 contain Tables 1--4; physical
page 21 contains Table 5. The appendices below retain each numerical
claim separately from any stronger conclusion of its companion proof.
An infinity symbol in the source's upper-bound column can mean that no
finite example was reported. It is not, by itself, a nonexistence theorem.

\begin{theorem}[Numerical synthesis]\label{thm:numerical}
Subject to the external premises explicitly specified by the cited
companion proofs, every lower inequality and every finite upper
inequality in Appendix~\ref{app:bounds} holds. Its 47 displayed exact
counts are counts of isomorphism classes at the stated upper order.
The 118 corrected symmetry minima in Appendix~\ref{app:symmetry} hold
in their stated symmetry classes.
\end{theorem}
Here ``exact'' applies only where the source asserts an exact count.
A supplied graph for a row marked ``at least one'' is not an exhaustion
of that row's extremal graphs. In particular, the last quartic row proved
here is a pair of inequalities, not an equality or an exact count.

\begin{proof}
The lower and exact-count columns refer to the numbered companion
proofs in Appendix~\ref{app:components}. Each of the 142 rows is joined
to its lower-bound proof; each of the 47 exact counts is separately
joined to its exhaustion or structural classification. The common
positive-witness proof C1 checks the 84 finite upper inequalities.
The symmetry proof C4 supplies the separate catalogue exhaustion and
symmetry certificates for the 118 rows. The only previously missing
lower inequality, for $(k,g,\lambda)=(4,6,1)$, follows from
Theorem~\ref{thm:57} below. The machine-readable concordance preserves
the literal source rows and identifies the exact companion editions;
the mathematical statements of those companions, rather than a numerical
agreement or a solver status, supply these implications.
\end{proof}

Among the source's highlighted conclusions are
\[
\begin{gathered}
n(3,6,2)=24,\qquad n(3,8,8)=40,\qquad
n(3,9,6)=n(3,9,8)=60,\\
n(4,5,1)=30,\qquad n(4,6,9)=35,\qquad n(6,5,20)=42.
\end{gathered}
\]
The first two use the cubic catalogue argument C9; the two order-60
statements use the unrestricted girth-nine argument C3, including its
explicit cage-classification premise. They do not follow from a
vertex-transitive minimum alone. The remaining three use C6, C30
together with the lower and positive proofs, and C2 with C1,
respectively. There are two isomorphism classes at $(24,3,6,2)$ and one
at $(40,3,8,8)$, $(30,4,5,1)$ and $(35,4,6,9)$. No exact count at order
60 or for the displayed order-42 row is asserted here.

\section{Cycle counting and complete generation}\label{sec:algorithm}
We state the mathematical interface of the complete-generation proof
\cite{algorithm}. It is independent of the historical running times or
the adequacy of fixed arrays in an old implementation.

\begin{lemma}[Deleted-edge counting]\label{lem:count}
Let $H$ have girth at least $g$ and let $uv$ be an edge. In $H-uv$,
perform breadth-first search from $u$. Set $a(u)=1$ and
\[
a(x)=\sum_{y\sim x,\ d(u,y)=d(u,x)-1}a(y).
\]
If $d(u,v)=g-1$, the number of $g$-cycles through $uv$ is $a(v)$;
otherwise it is zero. The computation takes $O(|V(H)|+|E(H)|)$
arithmetic operations.
\end{lemma}
\begin{proof}
Deleting $uv$ identifies its cycles with simple $u,v$ paths. Girth
forbids a path shorter than $g-1$, so a path of that length is shortest.
Every shortest path has a unique penultimate vertex in the preceding
layer, giving the recurrence. Before the FIFO queue removes a vertex,
all preceding-layer contributions have arrived. Each adjacency is
examined at most twice.
\end{proof}
For comparisons with $\lambda$, counts can be capped at $\lambda+1$.
The identity $\min(C,a+b)=\min(C,\min(C,a)+\min(C,b))$ proves that this
preserves the decisions below, equal to, or above $\lambda$. Arithmetic
complexity is not an assumption of constant-time unbounded integers;
the capped additions use $O(\log(\lambda+2))$ bits.

For odd girth $2s+1$, start with the degree-$k$ tree of depth $s$
about a vertex, grouping leaves by their first root edge. For even girth
$2s$, start with an edge and the two degree-$k$ trees of depth $s-1$
on its sides. Add isolates to reach the target order. Girth prevents
collisions inside these trees. A horizontal edge joins different leaf
groups. The tree path it closes is a girth cycle, and conversely every
girth cycle through a prescribed root edge (or the central edge in the
even case) arises this way. Consequently each group has horizontal
incidence quota $\lambda$.

\begin{theorem}[Corrected generation]\label{thm:generation}
The following serial depth-first scheme covers every target up to
isomorphism and, with exact terminal keys, outputs it once. Add only
edges between deficient vertices that preserve girth and do not force
more than $\lambda$ girth cycles through an edge. In the first phase
choose a \emph{deficient leaf group} and attempt all eligible horizontal
edges incident with that group. Once all quotas are met, choose an
unsaturated nonisolated vertex and attempt all eligible incident edges.
The scheme may use exact uncoloured keys on the entire graph, literal
rollback, historical sibling exclusions, and the least-labelled
isolated-neighbour representative as specified in \cite{algorithm}.
The final test requires connectedness, degree $k$ and cycle count
$\lambda$ at every edge.
\end{theorem}
\begin{proof}
Every target contains a seed. If a group is deficient, some missing
target edge is horizontal and incident with that group. In the second
phase an unsaturated nonisolated vertex likewise has a missing target
edge. If every nonisolated vertex is saturated while isolates remain,
there is no connected completion. Degree-capacity and cycle-excess tests
are necessary restrictions. Girth-cycle counts cannot decrease when
edges are added.

The less immediate pruning argument is chronological. A recursive call
strictly increases edge count, so a previous state with the same exact
uncoloured key cannot be an ancestor: its call has already finished.
Earlier sibling edges are excluded only after their entire child calls
finish. Inductively, a target completion removed by such a historical
exclusion is already represented by an earlier completed call. If no
target edge is historically excluded, the deficient-group or selected-
vertex argument supplies an eligible target edge. Relabelling isolated
vertices justifies retaining the least isolated neighbour; when its
representative edge is historically excluded, the same chronological
argument supplies earlier coverage. It is not necessary to assert that
the historical exclusion masks are isomorphic.

The complete induction, including capacity pruning in the presence of
these masks, is given in \cite{algorithm}. Restoration returns exactly
the parent graph and records before a new sibling exclusion is installed.
There are at most $nk/2$ edge additions on a branch. The exact terminal
test proves soundness, and the terminal isomorphism keys give uniqueness.
\end{proof}
The proof also admits the one-phase version and delayed excess-cycle
tests, with the same exact terminal predicate. A bounded memo table may
omit storing a nonterminal state; doing so loses a pruning opportunity,
not a completion. Duplicate-free output still requires the stated
terminal-key rule.

\subsection{The printed first-phase rule}
Printed Algorithm 2 chooses a minimum-option leaf, which can lie in a
group whose quota is already full. For $(n,k,g,\lambda)=(9,4,3,1)$,
take the star with centre 0, leaves $1,2,3,4$, and four isolates. After
adding $12$, leaves 1 and 2 each have four valid neighbours, all isolated;
leaves 3 and 4 have five, including each other. The printed rule chooses
1 or 2 but then has no horizontal edge to attempt. It misses $34$ and
therefore misses the $3\times3$ rook graph, which is connected,
4-regular and has exactly one triangle through each edge. The separate
literal counterexample checks all tie and first-edge histories.

The correction is to loop over a deficient group, as in
Theorem~\ref{thm:generation}. The later December 2024 C source, at
commit \texttt{41022f4a3ad84e4502f8889c5b64132225411ca2}, already uses
this corrected loop. The counterexample to the printed June rule is
not a counterexample to that later implementation. The abstract theorem
and exact reference implementation do not retrospectively certify every
historical machine execution; the numerical claims have their own proofs
and certificates.

\section{Quartic graphs with one hexagon through every edge}
Let $Q$ denote the hypotheses of a connected quartic graph of girth six
with exactly one hexagon through each edge. The following argument has
no bipartiteness assumption.

\subsection{The smaller-order reduction}
The two hexagons through a vertex pair its four incident edges. Distinct
hexagons meet in at most one vertex. The distance spheres of radii
$0,1,2,3$ about a vertex have sizes
\[
                         1,\quad4,\quad12,\quad34.
\]
Indeed there are 36 length-three nonbacktracking paths starting at the
vertex; precisely the two antipodal pairs on its two hexagons identify
their ends. A further collision would force an edge into two hexagons.
Thus $|V(G)|\ge51$. Counting edge--hexagon incidences gives
$6h=2|V(G)|$, so $3$ divides $|V(G)|$.

We import the complete order-51 spectral lemma \cite{reduction}.
At equality the radius-three ball is the whole graph, and the adjacency
matrix $A$ and the antipodal two-factor matrix $M$ satisfy
\[
             M=A^3+A^2-6A-3I-J.
\]
The proof compares the rational multiplicities of the cycle eigenvalues
of $M$ with the cubic preimages under this polynomial. Its cyclotomic
integer-ring premise and finite residue certificates are explicit in
\cite{reduction}. The resulting multiplicity congruences contradict the
principal eigenspace. This excludes order 51; we do not replace that
lemma with a numerical spectral approximation.

\subsection{A complete local cover at order 54}
For a hexagon $C$ of a hypothetical order-54 graph, write $t(C)$ for the
number of vertices at distance at least three from $C$. The local
incidence lemmas \cite{order54} give twelve first-layer vertices.
The second layer consists of vertices $S$ with one first-layer parent
and vertices $T$ with two parents. The remaining vertices form $Z$, with
\[
       |T|=|Z|=t(C),\qquad |S|=36-2t(C).
\]
The previously proved three- and four-exterior exclusions, together with
the layer restrictions, leave $t(C)\in\{0,5,6\}$.

For $t=5$, the auxiliary graph on $Z$ has the three coloured hosts and
actual parent assignments of the complete five-exterior theorem. Its
three necessary formulas have separate reconstructed refutations. Thus
$t=5$ is impossible. This uses the full five-exterior cover; no
four-exterior filter is transferred to it without a proof.

For $t=6$, each $T$ gives an edge between its two $Z$ neighbours. The
resulting simple graph $R$ has six vertices and six edges, with colours
recording the opposite anchors on $C$. The parent-cover theorem excludes
the fork and ear configurations and gives eight uncoloured shapes,
twelve coloured rows and seventy raw parent records. The letters
$N,L,R$ name actual parent choices. Disconnected hosts, a leaf, or an
isolated $Z$ vertex are allowed where the theorem permits them.

Each raw record gives a host with 54 vertices, 66 fixed edges and the
complete allowed optional $SS$, $SZ$ and $ZZ$ pairs. The fixed edges are
the six cycle edges, twelve edges to the first layer, twelve parent edges
to $T$, twenty-four parent edges to $S$, and twelve $TZ$ edges. Shared
parents change the individual deficits, not these totals. Optional
pairs include the isolated-$Z$ possibilities; their host degree can be
29. No degree-24 assumption from the five-exterior case is used.

Literal permutations of all 54 vertices take these seventy domains to
the twenty-four domains in Appendix~\ref{app:hosts}. Both the entire
fixed-edge set and the entire optional-edge set are checked under each
permutation. These equalities transport every graph completion and its
cycle properties. A claim that the representatives are mutually
nonisomorphic is unnecessary.

\subsection{Necessary formulas and their refutations}
There is one Boolean variable for each optional edge. The original
completion formula imposes degree four and the necessary girth
restrictions. Every one of the twelve fixed $TZ$ edges must occur in
a hexagon. Enumerating simple length-five paths between its ends in the
literal finite host gives the complete candidate list. Cycles are
discarded only for the proved degree or shortcut obstructions. A shared
positive gate is equivalent to the conjunction of the optional edges
of a cycle; the disjunction of these gates is required for each fixed
$TZ$ edge. The catalogue is therefore complete even when a host has an
isolated $Z$ vertex.

A short cycle, or two distinct hexagons sharing an edge, gives a
necessary negative clause on its optional edges. Every universal clause
and every subsequent lazy cut carries such a literal witness. Hence a
$Q$ completion would satisfy the entire final formula, after extending
its edge assignment to the gate variables. Only this necessary direction
is used.

The independent saved-certificate procedure recomputes the selected
catalogue, validates every cut witness, rebuilds the complete DIMACS
formula in its fixed clause order, and compares it with the formula
actually refuted. An unchanged reverse-unit-propagation checker then
validates all proof additions and an explicit empty clause. Its active
and historical clauses are consequences of that same original formula;
the solver's own UNSAT declaration is not a premise. Dimensions and
checked additions for all twenty-four domains appear in
Appendix~\ref{app:hosts}.

One domain's accepted proof is a separate copy of its original trace
followed by an empty-clause line. This is not an unconditional rule for
repairing a trace: the unchanged checker accepted the resulting complete
trace against the reconstructed original formula. A previous failed
read remains failed. Another domain uses a separately completed proof
of its already fixed formula. Neither case repeats or assumes the
correctness of an incomplete search.

It follows that no hexagon of an order-54 $Q$ graph can have $t=6$:
its completion would transport to one of the refuted domains and induce
a satisfying assignment of that domain's necessary formula.

\subsection{The global contradiction and lower bound}
\begin{theorem}\label{thm:54}
There is no order-54 graph satisfying $Q$.
\end{theorem}
\begin{proof}
There are 108 edges and therefore 18 hexagons. Their incidence graph
is simple and 6-regular: the two hexagons through each vertex are
distinct, and two hexagons meet at most once. The zero-neighbour lemma
\cite{order54} proves that two zero hexagons cannot meet a zero hexagon
at opposite vertices. Thus a zero hexagon has at most three zero
neighbours and at least three positive neighbours, where positive means
$t(C)>0$.

If $n_0,n_+$ count zero and positive hexagons, then
\[
                   3n_0\le6n_+,\qquad n_0+n_+=18,
\]
so $n_+\ge6$. But every hexagon has $t\in\{0,5,6\}$, and the complete
local refutations exclude both positive values. This is a contradiction.
The argument excludes the all-zero case; it does not assert that a
single zero hexagon is impossible independently of the others.
\end{proof}

\begin{theorem}\label{thm:57}
One has $57\le n(4,6,1)\le84$.
\end{theorem}
\begin{proof}
The sphere bound gives order at least 51 and the incidence count gives
divisibility by three. The order-51 lemma and Theorem~\ref{thm:54}
exclude the only possible orders below 57. The explicit 84-vertex graph
\texttt{UntilOrder250\_427} satisfies the defining properties by the
independent positive-witness tests C1. Thus the upper inequality holds.
\end{proof}
This proves exactly the lower/upper claim in source Table 2. It neither
excludes order 57 nor determines the exact minimum or number of graphs
at order 84.

The source's even-girth nonbipartite example also follows from the
separately classified $\egr(15,6,4,16)$ graph. In any nonempty regular
bipartite graph, counting edges from the two parts gives equal part
sizes. Its order is therefore even. The accepted order-15 witness is
nonbipartite, so the proposed bipartiteness analogue for minimum
edge-girth-regular graphs fails, as observed by the source authors.

\section{Literal certificates for the illustrated graphs}\label{sec:figures}
The following certificates complete the descriptive claims in the
source's Figures 2 and 4. They use the same three graphs as the positive
witness collection C1: \texttt{UntilOrder250\_50},
\texttt{UntilOrder250\_82} and \texttt{UntilOrder250\_300}, of orders
16, 20 and 40, respectively. All labels below are zero-based labels of
those exact graph6 records. An array $p$ lists $p(0),p(1),\ldots$;
line breaks within an array do not restart its indices. Complete arrays
and graph6 strings are also retained in the literal figure certificate.

\subsection{The Moebius--Kantor graph and its extension}
Define $G(8,3)$ on $u_i,v_i$, with $i$ taken modulo eight, by the edges
$u_iu_{i+1}$, $u_iv_i$ and $v_iv_{i+3}$. Label $u_i$ by $i$ and $v_i$
by $8+i$. The following full adjacency isomorphism identifies the
source's order-16 graph with this Moebius--Kantor graph:
\[
p=\bigl[\begin{gathered}
0,1,7,8,4,3,12,5,10,14,2,9,6,11,13,15
\end{gathered}\bigr].
\]
The permutations $u_i\mapsto u_{i+1}$, $v_i\mapsto v_{i+1}$
and $u_i\mapsto v_{3i}$, $v_i\mapsto u_{3i}$ preserve all three edge
types and act transitively together. Transporting them to the source
labels gives the two checked automorphisms
\[
a_1=\bigl[\begin{gathered}
1,10,0,11,7,4,14,12,13,15,5,8,2,6,9,3
\end{gathered}\bigr].
\]
\[
a_2=\bigl[\begin{gathered}
3,13,14,0,6,11,4,15,12,10,9,5,8,1,2,7
\end{gathered}\bigr].
\]
Their orbit of vertex zero contains all sixteen vertices. The following
injection identifies an induced copy of this graph in the source's
order-20 graph:
\[
\iota=\bigl[\begin{gathered}
1,8,17,16,2,19,18,11,3,4,9,10,14,12,15,13
\end{gathered}\bigr].
\]
The complementary set is $D=\{0,5,6,7\}$. It is independent;
every vertex outside $D$ has exactly one neighbour in $D$, and each
vertex of $D$ has four neighbours outside it. Checking all induced
adjacencies and these attachments verifies the whole add-four-vertices
construction in Figure 2. It also proves connectedness, since the
induced $G(8,3)$ is connected and each added vertex meets it.

Let $A$ be the adjacency matrix of this order-20 graph. Exact integer
multiplication gives
\[
 (A^6)_{vv}=\begin{cases}292,&v\in D,\\298,&v\notin D.\end{cases}
\]
This diagonal counts closed walks of length six and is invariant under
every automorphism. The two values therefore exclude vertex transitivity.
The literal cycle $(0,1,8,9,3)$ is a pentagon, so the graph is not bipartite.
To exclude edge transitivity, suppose a connected graph is edge-transitive.
The orbits of the two endpoints of one edge cover all vertices. If those
orbits coincide the graph is vertex-transitive. If they differ, every edge
joins the two orbits, giving a bipartition. Both alternatives fail here.
Thus the order-20 graph is neither vertex-transitive nor edge-transitive.

\subsection{The order-40 graph}
For the source's order-40 graph, the following three permutations are
checked on the entire adjacency relation:
\[
b_1=\bigl[\begin{gathered}
0,1,3,2,5,4,7,6,8,9,11,10,12,15,14,13,17,16,21,20,\\
19,18,23,22,33,31,30,32,29,28,26,25,27,24,34,38,39,37,35,36
\end{gathered}\bigr].
\]
\[
b_2=\bigl[\begin{gathered}
0,2,3,1,5,4,8,6,7,9,12,13,11,14,15,10,17,18,19,20,\\
21,16,24,25,26,27,28,29,32,30,33,22,31,23,35,36,37,39,34,38
\end{gathered}\bigr].
\]
\[
b_3=\bigl[\begin{gathered}
1,34,37,0,18,21,9,12,14,8,6,33,32,7,27,24,30,5,31,26,\\
4,25,10,16,20,15,38,36,11,17,13,35,19,39,22,29,3,23,2,28
\end{gathered}\bigr].
\]
Their generated orbit of vertex zero has forty vertices. This proves
the vertex-transitivity assertion of Figure 4. These are positive
automorphism and orbit certificates; no assertion that the displayed
generators generate the full automorphism group is needed.

All finite assertions in this section were independently checked from
the complete literal arrays, including the full closed-walk diagonal.
The separate companion arguments still supply the defining graph
parameters, uniqueness and extremality claims. The present certificates
do not infer those conclusions from symmetry or an embedding.


\section{External premises, corrections and attribution}
The companion proofs distinguish a direct certificate from a borrowed
classification. In particular, C9 uses the specified complete cubic
catalogues, C3 uses the classified cubic girth-nine cages, and C4 uses
the Poto\v cnik--Spiga--Verret catalogues of cubic vertex-transitive
graphs through order 1280 and tetravalent arc-transitive graphs through
order 640. Checking every retained record does not prove those external
catalogues complete. Positive symmetry certificates prove transitivity
by explicit automorphisms and orbits, without requiring a full
automorphism-group computation.

The elementary lower-bound and seed arguments retain the source's
attribution to Drglin, Filipovski, Jajcay and Raiman. The order-51 proof
retains its stated theorem about cyclotomic integer rings. Other borrowed
inputs and their exact applicability remain in the corresponding
companion manuscript; combining the rows does not discharge those
premises. The proof-source concordance includes each such manuscript
rather than replacing it with a claim of wholly independent rederivation.

There are two further corrections beyond the printed generation rule.
First, C4 proves
\[
                         n_{\vt}(3,14,56)=486,
\]
correcting 506 in Table 5. The literal catalogue graph has 729 edges
and 2916 fourteen-cycles, so $2916\cdot14=729\cdot56$; independent
edgewise tests and explicit vertex-transitivity witnesses supply its
positive certificate. The complete catalogue comparison excludes
smaller orders. The other 117 minima agree with the fixed source.
Second, this graph contradicts the unqualified physical-page-12 sentence
denying a cubic connected vertex-transitive graph on 486 vertices.
We do not silently insert a girth restriction into that sentence.

The source authors' refutations of the earlier extremality conjecture
are preserved as their results: the order-40 witness at $(3,8,8)$ is
smaller than 54, and the order-162 witness at $(3,12,32)$ is smaller
than 486. This synthesis makes no new priority claim for either.

The source's separate Conjecture 2 asserts four infinite minima. The
previously proved companion C34 establishes $n(3,8,14)=\infty$ by a
local forced-configuration and protected-edge argument independent of
the ambient order. This strengthens the source's finite lower bound 64
and resolves its fourth clause within the acquired component collection.
For the other three clauses this synthesis supplies only
\[
n(3,7,6)\ge42,\qquad n(3,8,10)\ge48,\qquad n(3,8,12)\ge48.
\]
It does not prove their proposed infinity values or assess their current
status elsewhere. This corrects the overbroad wording in an earlier
internal completion plan that described all four clauses as unresolved.
No new conjecture credit arises from assembling the already proved C34.

\section{Reproduction boundary}
The mathematical payload consists of the companion proofs, exact input
graphs and classifications with their stated provenance, local-domain
models and vertex permutations, formula construction and witness-checking
sources, and complete accepted refutations. A source-only rebuild of the
RUP checker is sufficient; a historical platform-specific executable is
not a mathematical premise. Formula construction, certificate verification
and finite coverage are distinct from rerunning certificate discovery.

All 49 component defaults and the five overarching route families
have separately accepted private executions. The overarching families
cover the 29 complete saved refutations, the seventy-to-twenty-four domain
transport, the literal figure certificates, the order-51 modular
certificates, and the two algorithm defaults. The current concordance
joins the 16 principal claim objects, all 260 source-table rows and all
53 direct written-source or composition records to those exact gates.
These counts do not remove the external premises stated above or grant
new conjecture or historical-running-time claims.

The staged scientific interfaces preserve the complete original defaults,
their controls and their full report comparisons. Their current editions
include the separately accepted catalogue-inventory and matching-manifest
repairs. The package-relative binding is
at \path{scripts/edge_girth_regular_graphs/execution/execute.py}; its
immutable plan declares 81 bounded application configurations. It binds
exact acquired worker sources to package paths and allows only explicitly
named package or environment fields to relocate. The literal figure
checker remains a separate relative entrypoint with its acquired
isolation boundary. The readable reproduction guide lists these routes
and distinguishes metadata preparation from actual scientific execution.

The publication closure retains the selected mathematical and
runtime payloads, all written-source joins, current adapters, mode and
runtime manifests, and full case expectations. Historical fixed-host
supervisors, an entire old receipt roster and old machine binaries are
not prerequisites merely because a prior gate records their identities.
Records actually read as scientific comparisons or written premises are
retained. Documentary identities remain explicit and are not silently
treated as newly executed files. No old artifact or Git history is
removed by this publication mechanism.

The first bounded binding targets the acquired macOS/Darwin ownership
mechanism, including its \texttt{libproc} process-group observations;
it is not a cross-platform process-reaping claim. Python and its required
packages, the compiler and sanitizer runtime, and the specified nauty
headers and library are explicit external environment premises. Canonical
and classification conventions remain those of the companion proofs.
The typed package-relative binding has a distinct source review. Its
metadata generator derives file identities and path projections from logical
roles, without requiring the author's old host paths as configuration keys.
The actual focused metadata controls checked all 81 case inverses and 20
rejections. A separately owned complete dense-bounds integration reproduced
its unchanged original report through this interface. These observations
cover the binding boundary and one scientific route in the recorded
environment. The publication admission at
\path{data/edge_girth_regular_graphs/admission.json} must separately bind
the exact assembly, manuscript, rendered PDF, registration and scoped
selected-checkout evidence. The accepted mathematical executions do not
by themselves assert that an arbitrary fresh machine can run every route.
The reproduction guide gives the concrete prerequisites and commands;
its instructions do not substitute for those actual admission records.

\appendix
\section{Bounds and asserted exact counts}\label{app:bounds}
$L_s$ is the fixed source's lower bound, $L$ the proved lower bound, and
$U_s$ its finite upper bound. A dash in $U_s$ means no finite upper
claim, not a proof of nonexistence. The column $c(U_s)$ gives only
asserted exact isomorphism counts; its dash leaves that count unspecified.
$C_L,C_c$ identify the lower and counting companion proofs below. All
finite upper witnesses use C1. An infinite $L$ is an all-order
nonexistence conclusion of its cited component.
{\small\subsection*{Source Table 1}
\begin{longtable}{rrrrrrrrl}
\toprule $k$ & $g$ & $\lambda$ & $L_s$ & $L$ & $U_s$ & $c(U_s)$ & $C_L$ & $C_c$ \\ \midrule\endhead
3 & 3 & 2 & 4 & 4 & 4 & 1 & 2 & 7 \\
3 & 4 & 2 & 8 & 8 & 8 & 1 & 2 & 7 \\
3 & 4 & 4 & 6 & 6 & 6 & 1 & 2 & 7 \\
3 & 5 & 2 & 20 & 20 & 20 & 1 & 2 & 9 \\
3 & 5 & 4 & 10 & 10 & 10 & 1 & 2 & 7 \\
3 & 6 & 2 & 24 & 24 & 24 & 2 & 9 & 9 \\
3 & 6 & 4 & 18 & 18 & 18 & 1 & 2 & 9 \\
3 & 6 & 6 & 16 & 16 & 16 & 1 & 2 & 9 \\
3 & 6 & 8 & 14 & 14 & 14 & 1 & 2 & 9 \\
3 & 7 & 2 & 42 & 42 & 56 & --- & 9 & --- \\
3 & 7 & 4 & 28 & 28 & 28 & 1 & 2 & 9 \\
3 & 7 & 6 & 42 & 42 & --- & --- & 9 & --- \\
3 & 7 & 8 & $\infty$ & $\infty$ & --- & --- & 2 & --- \\
3 & 8 & 2 & 56 & 56 & 64 & --- & 49 & --- \\
3 & 8 & 4 & 48 & 48 & 48 & --- & 9 & --- \\
3 & 8 & 6 & 48 & 48 & 48 & --- & 9 & --- \\
3 & 8 & 8 & 40 & 40 & 40 & 1 & 9 & 9 \\
3 & 8 & 10 & 48 & 48 & --- & --- & 9 & --- \\
3 & 8 & 12 & 48 & 48 & --- & --- & 9 & --- \\
3 & 8 & 14 & 64 & $\infty$ & --- & --- & 34 & --- \\
3 & 8 & 16 & 30 & 30 & 30 & 1 & 2 & 9 \\
\bottomrule\end{longtable}

\subsection*{Source Table 2}
\begin{longtable}{rrrrrrrrl}
\toprule $k$ & $g$ & $\lambda$ & $L_s$ & $L$ & $U_s$ & $c(U_s)$ & $C_L$ & $C_c$ \\ \midrule\endhead
4 & 3 & 1 & 9 & 9 & 9 & 1 & 2 & 7 \\
4 & 3 & 2 & 6 & 6 & 6 & 1 & 2 & 7 \\
4 & 3 & 3 & 5 & 5 & 5 & 1 & 2 & 7 \\
4 & 4 & 1 & 18 & 18 & 18 & 4 & 5 & 5 \\
4 & 4 & 2 & 13 & 13 & 13 & 1 & 2 & 7 \\
4 & 4 & 3 & 14 & 14 & 14 & 1 & 7 & 12 \\
4 & 4 & 4 & 21 & 21 & --- & --- & 23 & --- \\
4 & 4 & 5 & 10 & 10 & 10 & 1 & 2 & 7 \\
4 & 4 & 6 & 10 & 10 & 10 & 1 & 2 & 7 \\
4 & 4 & 7 & $\infty$ & $\infty$ & --- & --- & 2 & --- \\
4 & 4 & 8 & $\infty$ & $\infty$ & --- & --- & 2 & --- \\
4 & 4 & 9 & 8 & 8 & 8 & 1 & 2 & 7 \\
4 & 5 & 1 & 30 & 30 & 30 & 1 & 6 & 6 \\
4 & 5 & 2 & 30 & 30 & 30 & --- & 46 & --- \\
4 & 5 & 3 & 30 & 30 & 55 & --- & 46 & --- \\
4 & 5 & 4 & 30 & 30 & --- & --- & 46 & --- \\
4 & 5 & 5 & 24 & 24 & 24 & 1 & 45 & 45 \\
4 & 5 & 6 & 20 & 20 & 20 & 1 & 2 & 26 \\
4 & 5 & 7 & 35 & $\infty$ & --- & --- & 33 & --- \\
4 & 5 & 8 & 65 & $\infty$ & --- & --- & 32 & --- \\
4 & 5 & 9 & $\infty$ & $\infty$ & --- & --- & 2 & --- \\
4 & 6 & 1 & 57 & 57 & 84 & --- & 50 & --- \\
4 & 6 & 2 & 51 & 51 & 96 & --- & 31 & --- \\
4 & 6 & 3 & 45 & 47 & 90 & --- & 31 & --- \\
4 & 6 & 4 & 45 & 45 & 60 & --- & 31 & --- \\
4 & 6 & 5 & 42 & 45 & 81 & --- & 31 & --- \\
4 & 6 & 6 & 40 & 41 & 64 & --- & 31 & --- \\
4 & 6 & 7 & 39 & 39 & 60 & --- & 31 & --- \\
4 & 6 & 8 & 39 & 39 & 48 & --- & 31 & --- \\
4 & 6 & 9 & 35 & 35 & 35 & 1 & 2 & 30 \\
4 & 6 & 10 & 39 & 39 & 48 & --- & 44 & --- \\
4 & 6 & 11 & 36 & 36 & 42 & --- & 2 & --- \\
4 & 6 & 12 & 35 & 35 & 40 & --- & 43 & --- \\
4 & 6 & 13 & 36 & 36 & 60 & --- & 41 & --- \\
4 & 6 & 14 & 36 & 36 & --- & --- & 42 & --- \\
4 & 6 & 15 & 33 & 33 & 40 & --- & 40 & --- \\
4 & 6 & 16 & 33 & 33 & 36 & --- & 2 & --- \\
4 & 6 & 17 & 33 & 33 & --- & --- & 2 & --- \\
4 & 6 & 18 & 32 & 32 & 32 & 1 & 39 & 39 \\
4 & 6 & 19 & 36 & 36 & --- & --- & 38 & --- \\
4 & 6 & 20 & 36 & 36 & --- & --- & 37 & --- \\
4 & 6 & 21 & 30 & 30 & 30 & 1 & 25 & 29 \\
4 & 6 & 22 & 39 & $\infty$ & --- & --- & 36 & --- \\
4 & 6 & 23 & 42 & $\infty$ & --- & --- & 35 & --- \\
4 & 6 & 24 & 28 & 28 & 28 & 1 & 2 & 24 \\
4 & 6 & 25 & $\infty$ & $\infty$ & --- & --- & 2 & --- \\
4 & 6 & 26 & $\infty$ & $\infty$ & --- & --- & 2 & --- \\
4 & 6 & 27 & 26 & 26 & 26 & 1 & 2 & 11 \\
\bottomrule\end{longtable}

\subsection*{Source Table 3}
\begin{longtable}{rrrrrrrrl}
\toprule $k$ & $g$ & $\lambda$ & $L_s$ & $L$ & $U_s$ & $c(U_s)$ & $C_L$ & $C_c$ \\ \midrule\endhead
5 & 3 & 2 & 12 & 12 & 12 & 1 & 2 & 7 \\
5 & 3 & 4 & 6 & 6 & 6 & 1 & 2 & 7 \\
5 & 4 & 2 & 24 & 24 & 32 & --- & 2 & --- \\
5 & 4 & 4 & 16 & 16 & 16 & 1 & 2 & 13 \\
5 & 4 & 6 & 20 & 20 & --- & --- & 22 & --- \\
5 & 4 & 8 & 20 & 20 & --- & --- & 22 & --- \\
5 & 4 & 10 & 20 & 20 & --- & --- & 15 & --- \\
5 & 4 & 12 & 12 & 12 & 12 & 1 & 2 & 7 \\
5 & 4 & 14 & $\infty$ & $\infty$ & --- & --- & 2 & --- \\
5 & 4 & 16 & 10 & 10 & 10 & 1 & 2 & 7 \\
5 & 5 & 2 & 44 & 44 & 60 & --- & 48 & --- \\
5 & 5 & 4 & 38 & 38 & 66 & --- & 2 & --- \\
5 & 5 & 6 & 36 & 36 & --- & --- & 2 & --- \\
5 & 5 & 8 & 34 & 34 & 36 & --- & 2 & --- \\
5 & 5 & 10 & 32 & 32 & --- & --- & 2 & --- \\
5 & 5 & 12 & 32 & 32 & 32 & --- & 2 & --- \\
5 & 5 & 14 & 42 & 48 & --- & --- & 47 & --- \\
5 & 5 & 16 & $\infty$ & $\infty$ & --- & --- & 2 & --- \\
\bottomrule\end{longtable}

\subsection*{Source Table 4}
\begin{longtable}{rrrrrrrrl}
\toprule $k$ & $g$ & $\lambda$ & $L_s$ & $L$ & $U_s$ & $c(U_s)$ & $C_L$ & $C_c$ \\ \midrule\endhead
6 & 3 & 1 & 15 & 15 & 15 & 1 & 2 & 10 \\
6 & 3 & 2 & 12 & 12 & 12 & 1 & 2 & 7 \\
6 & 3 & 3 & 9 & 9 & 9 & 1 & 2 & 7 \\
6 & 3 & 4 & 8 & 8 & 8 & 1 & 2 & 7 \\
6 & 3 & 5 & 7 & 7 & 7 & 1 & 2 & 7 \\
6 & 4 & 1 & 32 & 36 & 40 & --- & 2 & --- \\
6 & 4 & 2 & 30 & 32 & 36 & --- & 2 & --- \\
6 & 4 & 3 & 28 & 28 & 32 & --- & 2 & --- \\
6 & 4 & 4 & 25 & 25 & 27 & --- & 2 & --- \\
6 & 4 & 5 & 24 & 24 & 24 & --- & 2 & --- \\
6 & 4 & 6 & 24 & 24 & 24 & --- & 21 & --- \\
6 & 4 & 7 & 24 & 24 & 36 & --- & 19 & --- \\
6 & 4 & 8 & 22 & 22 & 26 & --- & 20 & --- \\
6 & 4 & 9 & 20 & 20 & 20 & 2 & 2 & 19 \\
6 & 4 & 10 & 20 & 20 & 20 & 5 & 17 & 19 \\
6 & 4 & 11 & 20 & 20 & --- & --- & 2 & --- \\
6 & 4 & 12 & 20 & 20 & 20 & --- & 18 & --- \\
6 & 4 & 13 & 20 & 20 & 20 & --- & 16 & --- \\
6 & 4 & 14 & 18 & 18 & 18 & 1 & 16 & 17 \\
6 & 4 & 15 & 20 & 20 & --- & --- & 16 & --- \\
6 & 4 & 16 & 15 & 15 & 15 & 1 & 2 & 8 \\
6 & 4 & 17 & 16 & 16 & 16 & 1 & 2 & 14 \\
6 & 4 & 18 & 22 & 22 & --- & --- & 15 & --- \\
6 & 4 & 19 & 24 & 24 & --- & --- & 15 & --- \\
6 & 4 & 20 & 14 & 14 & 14 & 1 & 2 & 8 \\
6 & 4 & 21 & $\infty$ & $\infty$ & --- & --- & 2 & --- \\
6 & 4 & 22 & $\infty$ & $\infty$ & --- & --- & 2 & --- \\
6 & 4 & 23 & $\infty$ & $\infty$ & --- & --- & 2 & --- \\
6 & 4 & 24 & $\infty$ & $\infty$ & --- & --- & 2 & --- \\
6 & 4 & 25 & 12 & 12 & 12 & 1 & 2 & 7 \\
6 & 5 & 1 & 65 & 65 & --- & --- & 2 & --- \\
6 & 5 & 2 & 60 & 60 & --- & --- & 2 & --- \\
6 & 5 & 3 & 60 & 60 & --- & --- & 2 & --- \\
6 & 5 & 4 & 60 & 60 & --- & --- & 2 & --- \\
6 & 5 & 5 & 57 & 57 & --- & --- & 2 & --- \\
6 & 5 & 6 & 60 & 60 & --- & --- & 2 & --- \\
6 & 5 & 7 & 55 & 55 & --- & --- & 2 & --- \\
6 & 5 & 8 & 55 & 55 & --- & --- & 2 & --- \\
6 & 5 & 9 & 55 & 55 & --- & --- & 2 & --- \\
6 & 5 & 10 & 52 & 52 & --- & --- & 2 & --- \\
6 & 5 & 11 & 55 & 55 & --- & --- & 2 & --- \\
6 & 5 & 12 & 50 & 50 & --- & --- & 2 & --- \\
6 & 5 & 13 & 50 & 50 & --- & --- & 2 & --- \\
6 & 5 & 14 & 50 & 50 & --- & --- & 2 & --- \\
6 & 5 & 15 & 47 & 47 & 57 & --- & 2 & --- \\
6 & 5 & 16 & 50 & 50 & --- & --- & 2 & --- \\
6 & 5 & 17 & 45 & 45 & --- & --- & 2 & --- \\
6 & 5 & 18 & 45 & 45 & --- & --- & 2 & --- \\
6 & 5 & 19 & 45 & 45 & --- & --- & 2 & --- \\
6 & 5 & 20 & 42 & 42 & 42 & --- & 2 & --- \\
6 & 5 & 21 & 45 & 45 & --- & --- & 2 & --- \\
6 & 5 & 22 & 40 & 40 & 40 & 1 & 2 & 27 \\
6 & 5 & 23 & 45 & 45 & --- & --- & 28 & --- \\
6 & 5 & 24 & 45 & 45 & --- & --- & 28 & --- \\
6 & 5 & 25 & $\infty$ & $\infty$ & --- & --- & 2 & --- \\
\bottomrule\end{longtable}
}

\section{Corrected symmetry minima}\label{app:symmetry}
For $k=3$ the class is vertex-transitive, and for $k=4$ it is
arc-transitive. Both the printed value and the actual minimum are shown.
The complete comparison is C4.
{\small\begin{longtable}{rrrrr}
\toprule $k$ & $g$ & $\lambda$ & Source & Minimum \\ \midrule\endhead
3 & 9 & 2 & 108 & 108 \\
3 & 9 & 4 & 408 & 408 \\
3 & 9 & 6 & 60 & 60 \\
3 & 9 & 8 & 60 & 60 \\
3 & 10 & 2 & 160 & 160 \\
3 & 10 & 4 & 120 & 120 \\
3 & 10 & 10 & 112 & 112 \\
3 & 10 & 12 & 110 & 110 \\
3 & 10 & 16 & 90 & 90 \\
3 & 10 & 20 & 80 & 80 \\
3 & 11 & 2 & 1012 & 1012 \\
3 & 11 & 4 & 1012 & 1012 \\
3 & 11 & 8 & 506 & 506 \\
3 & 12 & 2 & 384 & 384 \\
3 & 12 & 4 & 384 & 384 \\
3 & 12 & 6 & 512 & 512 \\
3 & 12 & 8 & 272 & 272 \\
3 & 12 & 10 & 240 & 240 \\
3 & 12 & 12 & 256 & 256 \\
3 & 12 & 14 & 256 & 256 \\
3 & 12 & 16 & 234 & 234 \\
3 & 12 & 18 & 256 & 256 \\
3 & 12 & 20 & 234 & 234 \\
3 & 12 & 22 & 216 & 216 \\
3 & 12 & 24 & 204 & 204 \\
3 & 12 & 26 & 192 & 192 \\
3 & 12 & 28 & 182 & 182 \\
3 & 12 & 32 & 162 & 162 \\
3 & 12 & 34 & 168 & 168 \\
3 & 12 & 36 & 162 & 162 \\
3 & 13 & 26 & 384 & 384 \\
3 & 14 & 2 & 1092 & 1092 \\
3 & 14 & 6 & 1092 & 1092 \\
3 & 14 & 10 & 1008 & 1008 \\
3 & 14 & 14 & 768 & 768 \\
3 & 14 & 16 & 1008 & 1008 \\
3 & 14 & 28 & 504 & 504 \\
3 & 14 & 42 & 512 & 512 \\
3 & 14 & 44 & 406 & 406 \\
3 & 14 & 56 & 506 & 486 \\
3 & 15 & 10 & 1248 & 1248 \\
3 & 15 & 12 & 1280 & 1280 \\
3 & 15 & 20 & 864 & 864 \\
3 & 15 & 30 & 816 & 816 \\
3 & 15 & 32 & 620 & 620 \\
3 & 16 & 32 & 1250 & 1250 \\
3 & 16 & 72 & 1240 & 1240 \\
3 & 16 & 80 & 1012 & 1012 \\
3 & 16 & 96 & 1008 & 1008 \\
3 & 16 & 104 & 1280 & 1280 \\
4 & 7 & 1 & 252 & 252 \\
4 & 7 & 2 & 224 & 224 \\
4 & 7 & 3 & 273 & 273 \\
4 & 7 & 4 & 224 & 224 \\
4 & 7 & 7 & 160 & 160 \\
4 & 7 & 9 & 91 & 91 \\
4 & 7 & 14 & 80 & 80 \\
4 & 8 & 1 & 612 & 612 \\
4 & 8 & 2 & 320 & 320 \\
4 & 8 & 3 & 336 & 336 \\
4 & 8 & 4 & 252 & 252 \\
4 & 8 & 5 & 384 & 384 \\
4 & 8 & 6 & 288 & 288 \\
4 & 8 & 7 & 512 & 512 \\
4 & 8 & 8 & 243 & 243 \\
4 & 8 & 9 & 320 & 320 \\
4 & 8 & 10 & 240 & 240 \\
4 & 8 & 11 & 336 & 336 \\
4 & 8 & 12 & 288 & 288 \\
4 & 8 & 13 & 256 & 256 \\
4 & 8 & 14 & 256 & 256 \\
4 & 8 & 15 & 204 & 204 \\
4 & 8 & 16 & 240 & 240 \\
4 & 8 & 17 & 224 & 224 \\
4 & 8 & 18 & 256 & 256 \\
4 & 8 & 19 & 240 & 240 \\
4 & 8 & 20 & 162 & 162 \\
4 & 8 & 21 & 192 & 192 \\
4 & 8 & 22 & 256 & 256 \\
4 & 8 & 23 & 256 & 256 \\
4 & 8 & 24 & 126 & 126 \\
4 & 8 & 25 & 192 & 192 \\
4 & 8 & 26 & 160 & 160 \\
4 & 8 & 27 & 168 & 168 \\
4 & 8 & 28 & 144 & 144 \\
4 & 8 & 29 & 192 & 192 \\
4 & 8 & 30 & 160 & 160 \\
4 & 8 & 33 & 160 & 160 \\
4 & 8 & 36 & 140 & 140 \\
4 & 8 & 37 & 128 & 128 \\
4 & 8 & 39 & 128 & 128 \\
4 & 8 & 40 & 135 & 135 \\
4 & 8 & 41 & 128 & 128 \\
4 & 8 & 43 & 128 & 128 \\
4 & 8 & 45 & 128 & 128 \\
4 & 8 & 48 & 110 & 110 \\
4 & 8 & 60 & 100 & 100 \\
4 & 8 & 65 & 96 & 96 \\
4 & 9 & 6 & 504 & 504 \\
4 & 9 & 9 & 600 & 600 \\
4 & 9 & 15 & 546 & 546 \\
4 & 9 & 18 & 506 & 506 \\
4 & 9 & 19 & 504 & 504 \\
4 & 9 & 21 & 285 & 285 \\
4 & 9 & 27 & 320 & 320 \\
4 & 9 & 36 & 320 & 320 \\
4 & 10 & 50 & 625 & 625 \\
4 & 10 & 60 & 546 & 546 \\
4 & 10 & 61 & 640 & 640 \\
4 & 10 & 65 & 640 & 640 \\
4 & 10 & 67 & 640 & 640 \\
4 & 10 & 70 & 624 & 624 \\
4 & 10 & 71 & 640 & 640 \\
4 & 10 & 80 & 546 & 546 \\
4 & 10 & 90 & 512 & 512 \\
4 & 10 & 95 & 576 & 576 \\
4 & 10 & 108 & 420 & 420 \\
4 & 10 & 110 & 432 & 432 \\
\bottomrule\end{longtable}
}

\section{The twenty-four local refutations}\label{app:hosts}
``Raw domains'' counts the members of the complete seventy-record
parent cover transported to that representative. It does not count
graph isomorphism classes. Every row has a checked empty-clause
conclusion against its reconstructed necessary formula.
{\small\begin{longtable}{rp{.30\textwidth}rrrr}
\toprule ID & Shape and parents & Raw domains & Variables & Clauses & RUP additions \\ \midrule\endhead
0 & \texttt{c6\_opposite/NNN} & 1 & 15261 & 85921 & 445920 \\
1 & \texttt{c6\_opposite/NNL} & 6 & 15861 & 83147 & 43465 \\
4 & \texttt{c6\_opposite/NLL} & 6 & 16231 & 81194 & 21594 \\
5 & \texttt{c6\_opposite/NLR} & 6 & 16319 & 82798 & 1364 \\
13 & \texttt{c6\_opposite/LLL} & 6 & 16363 & 81451 & 1518 \\
16 & \texttt{c6\_opposite/LRL} & 2 & 16527 & 80934 & 211 \\
27 & \texttt{c6\_mixed/NNN} & 1 & 15197 & 84067 & 388681 \\
28 & \texttt{c6\_mixed/NLN} & 2 & 15797 & 82439 & 33160 \\
30 & \texttt{c5\_leaf\_far/NNN} & 1 & 14551 & 77874 & 7863 \\
31 & \texttt{c5\_leaf\_far/NNL} & 2 & 14917 & 76098 & 1212 \\
33 & \texttt{c5\_leaf\_near/NNN} & 1 & 14519 & 78032 & 5430 \\
34 & \texttt{square\_tail/NNN} & 1 & 15400 & 78367 & 1683 \\
35 & \texttt{square\_adj\_alternating/NNN} & 1 & 14344 & 73287 & 1327 \\
36 & \texttt{square\_adj\_three/NNN} & 1 & 14376 & 73415 & 1328 \\
37 & \texttt{square\_opp\_alternating/NNN} & 1 & 13880 & 75565 & 1315 \\
38 & \texttt{square\_opp\_alternating/NNL} & 2 & 14079 & 70981 & 1188 \\
40 & \texttt{square\_opp\_three/NNN} & 1 & 13880 & 73652 & 5143 \\
41 & \texttt{k23\_isolated/NNN} & 1 & 14114 & 72740 & 1069 \\
42 & \texttt{triangle\_tail/NNL} & 2 & 16095 & 80145 & 1403 \\
44 & \texttt{two\_triangles/NNL} & 6 & 17507 & 88782 & 13124 \\
47 & \texttt{two\_triangles/NLL} & 6 & 18118 & 91283 & 5217 \\
48 & \texttt{two\_triangles/NLR} & 6 & 18334 & 93832 & 3281 \\
56 & \texttt{two\_triangles/LLL} & 6 & 18567 & 92141 & 7115 \\
59 & \texttt{two\_triangles/LRL} & 2 & 18939 & 91689 & 1178 \\
\bottomrule\end{longtable}
}

\section{Companion proof concordance}\label{app:components}
The numbering in the numerical tables names the following exact
mathematical sources. Their byte identities and individual premise
statements are retained in the machine-readable source concordance.
The separate generation proof and order-51 reduction are also included
in that concordance.
\begin{longtable}{rp{.77\textwidth}}
\toprule ID & Companion mathematical proof \\ \midrule\endhead
1 & Cycle counts and explicit symmetry witnesses \\
2 & Lower bounds from short cycles \\
3 & Two minima among cubic graphs of girth nine \\
4 & Minimum orders within two symmetry classes \\
5 & Reconstruction from girth cycles \\
6 & A small incidence classification for unique pentagons \\
7 & Twenty extremal graphs of small order \\
8 & Two small extremal graphs from the triangle-free degree bound \\
9 & Cubic graphs from complete catalogues \\
10 & The graph of degree six with one triangle on each edge \\
11 & A graph at the girth-six Moore bound \\
12 & Four-cycle counts determine a fourteen-vertex graph \\
13 & Common-neighbour counts and the Clebsch graph \\
14 & A dense graph determined by its squares \\
15 & Squares near a shortest odd cycle \\
16 & Constant square counts on sixteen vertices \\
17 & A Latin square in the bipartite complement \\
18 & Excluding two odd orders by the same cycle partition \\
19 & Square counts on twenty vertices \\
20 & Eight squares through every edge \\
21 & Distance layers and square moments \\
22 & Square counts in graphs of degree five \\
23 & Four squares through each edge \\
24 & A graph from fourteen points and fourteen blocks \\
25 & Three vertices outside a Moore tree \\
26 & Three exterior vertices in a graph of girth five \\
27 & Reusing partial graphs in a classification of order forty \\
28 & Edges outside a Moore ball \\
29 & Four vertices outside a Moore edge tree \\
30 & Equality in the count of third neighbours \\
31 & Integer third-layer counts \\
32 & Eight pentagons through every edge \\
33 & Seven pentagons through every edge \\
34 & Fourteen octagons through every edge \\
35 & Twenty-three hexagons through every edge \\
36 & Twenty-two hexagons through every edge \\
37 & Twenty hexagons through every edge: a finite lower bound \\
38 & Nineteen hexagons through every edge: the remaining finite lower bound \\
39 & Eighteen hexagons through every edge: the unique graph of order thirty-two \\
40 & Fifteen hexagons through every edge: exclusion of order thirty-two \\
41 & Thirteen hexagons through every edge: exclusion of order thirty-three \\
42 & Fourteen hexagons through every edge: one exterior edge \\
43 & Twelve hexagons through every edge: one exterior edge \\
44 & Ten hexagons through every edge: three exterior edges \\
45 & Quartic graphs with five pentagons through every edge \\
46 & Quartic graphs with few pentagons through each edge \\
47 & A lower bound for pentavalent graphs with fourteen pentagons on each edge \\
48 & Two pentagons through every edge of a pentavalent graph \\
49 & Cubic graphs with two octagons through every edge \\
50 & Excluding order 54 with one hexagon through every edge \\
\bottomrule\end{longtable}


\begin{thebibliography}{99}
\bibitem{source} J. Goedgebeur and J. Jooken,
\emph{Exhaustive generation of edge-girth-regular graphs},
arXiv:2401.08271v2, 25 June 2024.
\bibitem{algorithm} Human Compression project,
\emph{Cycle counting and complete generation of edge-girth-regular graphs}.
Companion algorithm manuscript, including the corrected rule and exact
reference implementation.
\bibitem{reduction} Human Compression project,
\emph{A spectral obstruction and three remaining cases}. Companion order-51 spectral proof
and residue certificates.
\bibitem{order54} Human Compression project,
\emph{Excluding order 54 with one hexagon through every edge}.
Companion local proofs, complete domain transports and refutation
composition, C50.
\bibitem{C1} Human Compression project, \emph{Cycle counts and explicit symmetry witnesses}. Companion component C1; exact source and acquired premises are listed in the source concordance.
\bibitem{C2} Human Compression project, \emph{Lower bounds from short cycles}. Companion component C2; exact source and acquired premises are listed in the source concordance.
\bibitem{C3} Human Compression project, \emph{Two minima among cubic graphs of girth nine}. Companion component C3; exact source and acquired premises are listed in the source concordance.
\bibitem{C4} Human Compression project, \emph{Minimum orders within two symmetry classes}. Companion component C4; exact source and acquired premises are listed in the source concordance.
\bibitem{C5} Human Compression project, \emph{Reconstruction from girth cycles}. Companion component C5; exact source and acquired premises are listed in the source concordance.
\bibitem{C6} Human Compression project, \emph{A small incidence classification for unique pentagons}. Companion component C6; exact source and acquired premises are listed in the source concordance.
\bibitem{C7} Human Compression project, \emph{Twenty extremal graphs of small order}. Companion component C7; exact source and acquired premises are listed in the source concordance.
\bibitem{C8} Human Compression project, \emph{Two small extremal graphs from the triangle-free degree bound}. Companion component C8; exact source and acquired premises are listed in the source concordance.
\bibitem{C9} Human Compression project, \emph{Cubic graphs from complete catalogues}. Companion component C9; exact source and acquired premises are listed in the source concordance.
\bibitem{C10} Human Compression project, \emph{The graph of degree six with one triangle on each edge}. Companion component C10; exact source and acquired premises are listed in the source concordance.
\bibitem{C11} Human Compression project, \emph{A graph at the girth-six Moore bound}. Companion component C11; exact source and acquired premises are listed in the source concordance.
\bibitem{C12} Human Compression project, \emph{Four-cycle counts determine a fourteen-vertex graph}. Companion component C12; exact source and acquired premises are listed in the source concordance.
\bibitem{C13} Human Compression project, \emph{Common-neighbour counts and the Clebsch graph}. Companion component C13; exact source and acquired premises are listed in the source concordance.
\bibitem{C14} Human Compression project, \emph{A dense graph determined by its squares}. Companion component C14; exact source and acquired premises are listed in the source concordance.
\bibitem{C15} Human Compression project, \emph{Squares near a shortest odd cycle}. Companion component C15; exact source and acquired premises are listed in the source concordance.
\bibitem{C16} Human Compression project, \emph{Constant square counts on sixteen vertices}. Companion component C16; exact source and acquired premises are listed in the source concordance.
\bibitem{C17} Human Compression project, \emph{A Latin square in the bipartite complement}. Companion component C17; exact source and acquired premises are listed in the source concordance.
\bibitem{C18} Human Compression project, \emph{Excluding two odd orders by the same cycle partition}. Companion component C18; exact source and acquired premises are listed in the source concordance.
\bibitem{C19} Human Compression project, \emph{Square counts on twenty vertices}. Companion component C19; exact source and acquired premises are listed in the source concordance.
\bibitem{C20} Human Compression project, \emph{Eight squares through every edge}. Companion component C20; exact source and acquired premises are listed in the source concordance.
\bibitem{C21} Human Compression project, \emph{Distance layers and square moments}. Companion component C21; exact source and acquired premises are listed in the source concordance.
\bibitem{C22} Human Compression project, \emph{Square counts in graphs of degree five}. Companion component C22; exact source and acquired premises are listed in the source concordance.
\bibitem{C23} Human Compression project, \emph{Four squares through each edge}. Companion component C23; exact source and acquired premises are listed in the source concordance.
\bibitem{C24} Human Compression project, \emph{A graph from fourteen points and fourteen blocks}. Companion component C24; exact source and acquired premises are listed in the source concordance.
\bibitem{C25} Human Compression project, \emph{Three vertices outside a Moore tree}. Companion component C25; exact source and acquired premises are listed in the source concordance.
\bibitem{C26} Human Compression project, \emph{Three exterior vertices in a graph of girth five}. Companion component C26; exact source and acquired premises are listed in the source concordance.
\bibitem{C27} Human Compression project, \emph{Reusing partial graphs in a classification of order forty}. Companion component C27; exact source and acquired premises are listed in the source concordance.
\bibitem{C28} Human Compression project, \emph{Edges outside a Moore ball}. Companion component C28; exact source and acquired premises are listed in the source concordance.
\bibitem{C29} Human Compression project, \emph{Four vertices outside a Moore edge tree}. Companion component C29; exact source and acquired premises are listed in the source concordance.
\bibitem{C30} Human Compression project, \emph{Equality in the count of third neighbours}. Companion component C30; exact source and acquired premises are listed in the source concordance.
\bibitem{C31} Human Compression project, \emph{Integer third-layer counts}. Companion component C31; exact source and acquired premises are listed in the source concordance.
\bibitem{C32} Human Compression project, \emph{Eight pentagons through every edge}. Companion component C32; exact source and acquired premises are listed in the source concordance.
\bibitem{C33} Human Compression project, \emph{Seven pentagons through every edge}. Companion component C33; exact source and acquired premises are listed in the source concordance.
\bibitem{C34} Human Compression project, \emph{Fourteen octagons through every edge}. Companion component C34; exact source and acquired premises are listed in the source concordance.
\bibitem{C35} Human Compression project, \emph{Twenty-three hexagons through every edge}. Companion component C35; exact source and acquired premises are listed in the source concordance.
\bibitem{C36} Human Compression project, \emph{Twenty-two hexagons through every edge}. Companion component C36; exact source and acquired premises are listed in the source concordance.
\bibitem{C37} Human Compression project, \emph{Twenty hexagons through every edge: a finite lower bound}. Companion component C37; exact source and acquired premises are listed in the source concordance.
\bibitem{C38} Human Compression project, \emph{Nineteen hexagons through every edge: the remaining finite lower bound}. Companion component C38; exact source and acquired premises are listed in the source concordance.
\bibitem{C39} Human Compression project, \emph{Eighteen hexagons through every edge: the unique graph of order thirty-two}. Companion component C39; exact source and acquired premises are listed in the source concordance.
\bibitem{C40} Human Compression project, \emph{Fifteen hexagons through every edge: exclusion of order thirty-two}. Companion component C40; exact source and acquired premises are listed in the source concordance.
\bibitem{C41} Human Compression project, \emph{Thirteen hexagons through every edge: exclusion of order thirty-three}. Companion component C41; exact source and acquired premises are listed in the source concordance.
\bibitem{C42} Human Compression project, \emph{Fourteen hexagons through every edge: one exterior edge}. Companion component C42; exact source and acquired premises are listed in the source concordance.
\bibitem{C43} Human Compression project, \emph{Twelve hexagons through every edge: one exterior edge}. Companion component C43; exact source and acquired premises are listed in the source concordance.
\bibitem{C44} Human Compression project, \emph{Ten hexagons through every edge: three exterior edges}. Companion component C44; exact source and acquired premises are listed in the source concordance.
\bibitem{C45} Human Compression project, \emph{Quartic graphs with five pentagons through every edge}. Companion component C45; exact source and acquired premises are listed in the source concordance.
\bibitem{C46} Human Compression project, \emph{Quartic graphs with few pentagons through each edge}. Companion component C46; exact source and acquired premises are listed in the source concordance.
\bibitem{C47} Human Compression project, \emph{A lower bound for pentavalent graphs with fourteen pentagons on each edge}. Companion component C47; exact source and acquired premises are listed in the source concordance.
\bibitem{C48} Human Compression project, \emph{Two pentagons through every edge of a pentavalent graph}. Companion component C48; exact source and acquired premises are listed in the source concordance.
\bibitem{C49} Human Compression project, \emph{Cubic graphs with two octagons through every edge}. Companion component C49; exact source and acquired premises are listed in the source concordance.

\end{thebibliography}
\end{document}
